% Generated draft: requires mathematical and source review.
\documentclass[11pt]{article}
\usepackage[margin=25mm]{geometry}
\usepackage{fontspec}
\setmainfont{Latin Modern Roman}
\usepackage{amsmath,amssymb,mathtools,unicode-math}
\setmathfont{Latin Modern Math}
\usepackage{enumitem,needspace,xurl,hyperref}
\setlength{\parindent}{0pt}
\setlength{\parskip}{5pt}
\setlength{\emergencystretch}{3em}
\newcommand{\N}{\mathbb{N}}
\newcommand{\Z}{\mathbb{Z}}
\newcommand{\Q}{\mathbb{Q}}
\newcommand{\R}{\mathbb{R}}
\newcommand{\C}{\mathbb{C}}
\newcommand{\F}{\mathbb{F}}
\newcommand{\A}{\mathbb{A}}
\newcommand{\PP}{\mathbb{P}}
\newcommand{\E}{\mathbb{E}}
\newcommand{\eps}{\varepsilon}
\newcommand{\dd}{\,\mathrm{d}}
\begin{document}
\section*{UnsolvedMath ID: 800401}
\textbf{Source problem code:} catalogue entry
\subsection{Definitions and mathematical statement}
Let \(\bar K\) be an algebraically closed field of characteristic \(0\). The supplied source writes \(X/K\) without explicitly specifying \(K\) or its relationship to \(\bar K\). The formulation below uses the intended reading \(K\subseteq\bar K\), with \(Y=X\times_K\bar K\), and reads \(\bar K(X)\) as the rational-function field \(\bar K(Y)\) under the usual integral-variety convention. These base-field and variety conventions require confirmation from the source.

Let \(X/K\) be a quasi-projective variety and let \(f:X\to X\) be a dominant morphism, meaning that its image is Zariski dense. Write \(F=f_{\bar K}:Y\to Y\) for the induced morphism. For \(P\in Y(\bar K)=X(\bar K)\), define its forward orbit by
\[
\mathcal O_F(P)=\{F^n(P):n\in\Z_{\geq 0}\},\qquad F^0=\operatorname{id}_Y.
\]
A rational function \(g\in\bar K(Y)\) is invariant under \(F\) if \(g\circ F=g\) as rational functions; this equality is understood on their common domain of definition.

\textbf{Orbit Density Conjecture (supplied Conjecture 24).} Subject to the unresolved conventions above, if
\[
\text{there is no }g\in\bar K(Y)\setminus\bar K\text{ such that }g\circ F=g,
\]
then there exists \(P\in Y(\bar K)\) such that
\[
\overline{\mathcal O_F(P)}^{\mathrm{Zar}}=Y.
\]
The supplied record labels this conjecture open; that label is not an independently verified literature assessment.
\subsection{Short English statement}
The conjecture asserts that a dominant self-morphism of a quasi-projective variety in characteristic zero has a Zariski-dense forward orbit if it preserves no nonconstant rational function. The source's base-field notation and the conventions needed for its rational-function field require clarification.
\subsection{Sources}
\begin{itemize}
\item Local UnsolvedMath registry, permanent ID \(800401\), source problem code \texttt{catalogue entry}, titled \emph{Open Problems in Arithmetic Dynamics --- Conjecture 24 (Orbit Density Conjecture)}. The supplied status is \texttt{open}.
\item Benedetto, DeMarco, Ingram, Jones et al., \emph{Current Trends and Open Problems in Arithmetic Dynamics} (2018), Conjecture 24. Supplied source URL: \url{https://arxiv.org/abs/1806.04980}. The record reports access on 2026-10-16 and extraction by \texttt{arxiv\_latex}; these are supplied provenance details, not a claim of independent access or verification.
\item The supplied statement attributes the conjecture to Medvedev and Scanlon (2009), reference label \texttt{arxiv0901.2352}, Conjecture 5.10; and Amerik, Bogomolov and Rovinsky (2011), reference label \texttt{MR2862064}, Conjecture 1.2. No further bibliographic details for these references are supplied.
\end{itemize}
\end{document}
